\subsection{Testing section: CMS Components}\label{subsec:testing-section-cms-components}
The system is based on a role-based access control model, which enables the administrators to create users
and assign them to groups.
Users inherit privileges from the groups they are part of.
The framework defines several special groups as a part of initial configuration.
For example, the Admin group grants implicit access to the administration dashboard.
Even if a user are belongs a custom group,
that grants the privilege to Page editing,
without assigning the Admin group the user cannot log in to administration dashboard.

The system lacks proper descriptions for these special groups, and requires implicit understanding of the special groups.

\subsection{Testing section: AI Generation}\label{subsec:testing-section-ai-generation}
The framework provides couple of ways to generate content more efficiently by using AI\@.
The following table describes whether the feature was used and following paragraphs describe the quality of generated content.

\begin{table}[H]
    \centering
    \begin{tabular}{|p{0.4\linewidth}|c|c|c|}
        \hline
        \textbf{Feature} & \textbf{Corporate} & \textbf{Content-Driven} & \textbf{Personal} \\
        \hline
        Page Structure Generation & Yes & Yes & No \\
        \hline
        Article Generation & Yes & Yes & Yes \\
        \hline
        Static Page Generation & Yes & Yes & No \\
        \hline
        Page Component Generation & Yes & No & No \\
        \hline
        Documentation Generation & No & No & Yes \\
        \hline
    \end{tabular}
    \caption{Testing Section - AI Generation Overview}
    \label{tab:evaluation-test-ai-overview}
\end{table}

\subsubsection*{Page Structure Generation}
The generated structures are usable but a few pages need to be removed for each project.
For example, sometimes the Home page needs to be removed because it is implemented via the Home component.
A few times the AI generated the whole structure with Home being the parent.
Due to the system limitation not allowing transferring parent-child relation to other pages the whole structure
has to be deleted and regenerated.

In the content-driven scenario, the generated structure was ready to use and only blog articles had to be filled in manually.

\subsubsection*{Article Generation}
The generated output matched the desired content but the Markdown formatting had to be manually adjusted.
This is due to the use of a custom-made Markdown compiler, where some AI-Generated formatting choices
lead to undefined results.

\subsubsection*{Static Page Generation}
The static contact pages are generated using AI with simple prompt without styling guides
and detailed description of each scenario.
Combined with the use of a smaller model, this led to lower-quality results with mismatched theme styling
and incorrect typography.

\subsubsection*{Page Component Generation}
The page content generation is used to generate static table for corporate scenario.
With detailed prompt the output is visually presentable but not directly usable in production.
The issue is that generated components cannot expose editable properties or content to the editor.
If the user wants to edit the content, the change must be explicitly stated in the prompt rather than editing the output directly.

This makes the feature well-suited for drafting components in natural language and handing the output to development team
to implement fully functional component.

\subsubsection*{Documentation Generation}

The framework employs on-demand documentation generation, which shifts the responsibility of correctness validation from the system to the project maintainer.

This design introduces a trade-off between accessibility of undocumented code and reliability of the generated documentation.

\subsubsection*{Documentation Generation}
The Ai-Generated documentation outlines source files and lists methods of defined within each file.
The introduction is concise, the same applies to the method descriptions.
Due to the probabilistic nature of LLMs, the generated output varies in correctness, which is problematic in fact-driven environment,
such as technical documentation.
Additionally, the output is in Markdown format and inherits the same problems with formatting as Article Generation.

This means that for each documentation page the project maintainer needs to verify the correctness of the generated file,
and fix formatting before it can be declared production ready.
The issue is that the framework uses on-demand documentation generation,
which may expose users to incorrect and visually broken technical documentation.

This design decision introduces a trade-off between accessibility and reliability of the documentation.

\subsection{Testing section: Metrics Evaluation}\label{subsec:testing-section:-metrics-evaluation}
The framework is tested on the metrics defined in the table.
Some values are missing metrics about bug fixing the framework.
If critical bug is encountered during the deployment, the development of fixing the bug is not counted to the metrics.

\begin{table}[H]
    \centering
    \begin{tabular}{|p{0.23\linewidth}|p{0.23\linewidth}|p{0.23\linewidth}|p{0.23\linewidth}|}
        \hline
        \textbf{Evaluation Metrics}
        &
        \textbf{Corporate}
        &
        \textbf{Content-Driven}
        &
        \textbf{Personal} \\
        \hline
        Requirement Coverage
        &
        The only thing that had to be implemented by hand is the the homepage and since it is intended, the Requirement
        Coverage is 100\%
        &
        &
        \\
        \hline
        Number of Manual Interventions
        &
        Only one, homepage implementation
        &
        &
        \\
        \hline
        Time to Publishable Draft
        &
        Measuring the time between initial commit and final commit (not counting update commits) the deployment of the website
        took exactly one hour and two minutes.
        &
        &\\
        \hline
        Responsiveness
        &
        All the pages that are visible to the public are fully responsive although some elements are implemented
        in functional but not user-friendly way.
        The elements are header menu items on mobile devices.
        The administration websites require custom implementation of core user controls, and thus it is only optimized
        for desktop devices with standard landscape resolutions.
        &
        &\\
        \hline
        Basic SEO
        &
        Publicly accessible websites pass the SEO test if all metadata for a page are filled in.
        The only exception is the homepage which requires manual description to be added by editing the Home component
        directly.
        &
        &\\
        \hline
    \end{tabular}
    \caption{Testing Section - Metrics Evaluation}
    \label{tab:evaluation-metrics}
\end{table}
